\section{Functoriality of Proj}
\label{section-functoriality-proj}

\noindent
A graded ring map $\psi : A \to B$ does not always give rise to a morphism of
associated projective homogeneous spectra. The reason is that
the inverse image $\psi^{-1}(\mathfrak q)$
of a homogeneous prime $\mathfrak q \subset B$ may
contain the irrelevant prime $A_{+}$ even if $\mathfrak q$ does not
contain $B_{+}$.
The correct result is stated as follows.

\begin{lemma}
\label{lemma-morphism-proj}
Let $A$, $B$ be two graded rings.
Set $X = \text{Proj}(A)$ and $Y = \text{Proj}(B)$.
Let $\psi : A \to B$ be a graded ring map.
Set
$$
U(\psi)
=
\bigcup\nolimits_{f \in A_{+}\ \text{homogeneous}} D_{+}(\psi(f))
\subset Y.
$$
Then there is a canonical morphism of schemes
$$
r_\psi :
U(\psi)
\longrightarrow
X
$$
and a map of $\mathbf{Z}$-graded $\mathcal{O}_{U(\psi)}$-algebras
$$
\theta = \theta_\psi :
r_\psi^*\left(
\bigoplus\nolimits_{d \in \mathbf{Z}} \mathcal{O}_X(d)
\right)
\longrightarrow
\bigoplus\nolimits_{d \in \mathbf{Z}} \mathcal{O}_{U(\psi)}(d).
$$
The triple $(U(\psi), r_\psi, \theta)$ is
characterized by the following properties:
\begin{enumerate}
\item For every $d \geq 0$ the diagram
$$
\xymatrix{
A_d \ar[d] \ar[rr]_{\psi} & &
B_d \ar[d] \\
\Gamma(X, \mathcal{O}_X(d)) \ar[r]^-\theta &
\Gamma(U(\psi), \mathcal{O}_Y(d)) &
\Gamma(Y, \mathcal{O}_Y(d)) \ar[l]
}
$$
is commutative.
\item For any $f \in A_{+}$ homogeneous
we have $r_\psi^{-1}(D_{+}(f)) = D_{+}(\psi(f))$ and
the restriction of $r_\psi$ to $D_{+}(\psi(f))$
corresponds to the ring map
$A_{(f)} \to B_{(\psi(f))}$ induced by $\psi$.
\end{enumerate}
\end{lemma}

\begin{proof}
Clearly condition (2) uniquely determines the morphism of schemes
and the open subset $U(\psi)$. Pick $f \in A_d$ with $d \geq 1$.
Note that
$\mathcal{O}_X(n)|_{D_{+}(f)}$ corresponds to the
$A_{(f)}$-module $(A_f)_n$ and that
$\mathcal{O}_Y(n)|_{D_{+}(\psi(f))}$ corresponds to the
$B_{(\psi(f))}$-module $(B_{\psi(f)})_n$. In other words $\theta$
when restricted to $D_{+}(\psi(f))$ corresponds to a map of
$\mathbf{Z}$-graded $B_{(\psi(f))}$-algebras
$$
A_f \otimes_{A_{(f)}} B_{(\psi(f))}
\longrightarrow
B_{\psi(f)}
$$
Condition (1) determines the images of all elements of $A$.
Since $f$ is an invertible element which is mapped to $\psi(f)$
we see that $1/f^m$ is mapped to $1/\psi(f)^m$. It easily follows
from this that $\theta$ is uniquely determined, namely it is
given by the rule
$$
a/f^m \otimes b/\psi(f)^e \longmapsto \psi(a)b/\psi(f)^{m + e}.
$$
To show existence we remark that the proof of uniqueness above gave
a well defined prescription for the morphism $r$ and the map $\theta$
when restricted to every standard open of the form
$D_{+}(\psi(f)) \subset U(\psi)$ into $D_{+}(f)$.
Call these $r_f$ and $\theta_f$.
Hence we only need to verify that if $D_{+}(f) \subset D_{+}(g)$
for some $f, g \in A_{+}$ homogeneous, then the restriction of
$r_g$ to $D_{+}(\psi(f))$ matches $r_f$. This is clear from the
formulas given for $r$ and $\theta$ above.
\end{proof}

\begin{lemma}
\label{lemma-morphism-proj-transitive}
Let $A$, $B$, and $C$ be graded rings.
Set $X = \text{Proj}(A)$, $Y = \text{Proj}(B)$ and $Z = \text{Proj}(C)$.
Let $\varphi : A \to B$, $\psi : B \to C$ be graded ring maps.
Then we have
$$
U(\psi \circ \varphi) = r_\psi^{-1}(U(\varphi))
\quad
\text{and}
\quad
r_{\psi \circ \varphi}
=
r_\varphi \circ r_\psi|_{U(\psi \circ \varphi)}.
$$
In addition we have
$$
\theta_\psi \circ r_\psi^*\theta_\varphi
=
\theta_{\psi \circ \varphi}
$$
with obvious notation.
\end{lemma}

\begin{proof}
Omitted.
\end{proof}

\begin{lemma}
\label{lemma-surjective-graded-rings-map-proj}
With hypotheses and notation as in Lemma \ref{lemma-morphism-proj} above.
Assume $A_d \to B_d$ is surjective for all $d \gg 0$. Then
\begin{enumerate}
\item $U(\psi) = Y$,
\item $r_\psi : Y \to X$ is a closed immersion, and
\item the maps $\theta : r_\psi^*\mathcal{O}_X(n) \to \mathcal{O}_Y(n)$
are surjective but not isomorphisms in general (even if $A \to B$ is
surjective).
\end{enumerate}
\end{lemma}

\begin{proof}
Part (1) follows from the definition of $U(\psi)$ and the fact that
$D_{+}(f) = D_{+}(f^n)$ for any $n > 0$. For $f \in A_{+}$ homogeneous
we see that $A_{(f)} \to B_{(\psi(f))}$ is surjective because
any element of $B_{(\psi(f))}$ can be represented by a fraction
$b/\psi(f)^n$ with $n$ arbitrarily large (which forces the degree of
$b \in B$ to be large). This proves (2).
The same argument shows the map
$$
A_f \to B_{\psi(f)}
$$
is surjective which proves the surjectivity of $\theta$.
For an example where this map is not an isomorphism
consider the graded ring $A = k[x, y]$ where $k$ is a field
and $\deg(x) = 1$, $\deg(y) = 2$. Set $I = (x)$, so that
$B = k[y]$. Note that $\mathcal{O}_Y(1) = 0$ in this case.
But it is easy to see that $r_\psi^*\mathcal{O}_X(1)$
is not zero. (There are less silly examples.)
\end{proof}

\begin{lemma}
\label{lemma-eventual-iso-graded-rings-map-proj}
With hypotheses and notation as in Lemma \ref{lemma-morphism-proj} above.
Assume $A_d \to B_d$ is an isomorphism for all $d \gg 0$. Then
\begin{enumerate}
\item $U(\psi) = Y$,
\item $r_\psi : Y \to X$ is an isomorphism, and
\item the maps $\theta : r_\psi^*\mathcal{O}_X(n) \to \mathcal{O}_Y(n)$
are isomorphisms.
\end{enumerate}
\end{lemma}

\begin{proof}
We have (1) by Lemma \ref{lemma-surjective-graded-rings-map-proj}.
Let $f \in A_{+}$ be homogeneous. The assumption on $\psi$ implies that
$A_f \to B_f$ is an isomorphism (details omitted). Thus it is clear that
$r_\psi$ and $\theta$ restrict to isomorphisms over $D_{+}(f)$.
The lemma follows.
\end{proof}

\begin{lemma}
\label{lemma-surjective-graded-rings-generated-degree-1-map-proj}
With hypotheses and notation as in Lemma \ref{lemma-morphism-proj} above.
Assume $A_d \to B_d$ is surjective for $d \gg 0$ and that $A$ is generated
by $A_1$ over $A_0$. Then
\begin{enumerate}
\item $U(\psi) = Y$,
\item $r_\psi : Y \to X$ is a closed immersion, and
\item the maps $\theta : r_\psi^*\mathcal{O}_X(n) \to \mathcal{O}_Y(n)$
are isomorphisms.
\end{enumerate}
\end{lemma}

\begin{proof}
By Lemmas \ref{lemma-eventual-iso-graded-rings-map-proj} and
\ref{lemma-morphism-proj-transitive}
we may replace $B$ by the image of $A \to B$
without changing $X$ or the sheaves $\mathcal{O}_X(n)$.
Thus we may assume that $A \to B$ is surjective. By
Lemma \ref{lemma-surjective-graded-rings-map-proj} we get (1) and (2)
and surjectivity in (3).
By Lemma \ref{lemma-apply-modules} we see that both
$\mathcal{O}_X(n)$ and $\mathcal{O}_Y(n)$
are invertible. Hence $\theta$ is an isomorphism.
\end{proof}

\begin{lemma}
\label{lemma-base-change-map-proj}
\begin{reference}
\cite[II, Proposition 2.8.10]{EGA}
\end{reference}
With hypotheses and notation as in Lemma \ref{lemma-morphism-proj} above.
Assume there exists a ring map $R \to A_0$ and a ring map
$R \to R'$ such that $B = R' \otimes_R A$. Then
\begin{enumerate}
\item $U(\psi) = Y$,
\item the diagram
$$
\xymatrix{
Y = \text{Proj}(B) \ar[r]_{r_\psi} \ar[d] &
\text{Proj}(A) = X \ar[d] \\
\Spec(R') \ar[r] &
\Spec(R)
}
$$
is a fibre product square, and
\item the maps $\theta : r_\psi^*\mathcal{O}_X(n) \to \mathcal{O}_Y(n)$
are isomorphisms.
\end{enumerate}
\end{lemma}

\begin{proof}
This follows immediately by looking at what happens over the standard
opens $D_{+}(f)$ for $f \in A_{+}$.
\end{proof}

\begin{lemma}
\label{lemma-localization-map-proj}
With hypotheses and notation as in Lemma \ref{lemma-morphism-proj} above.
Assume there exists a $g \in A_0$ such that $\psi$ induces an
isomorphism $A_g \to B$. Then
$U(\psi) = Y$, $r_\psi : Y \to X$ is an open immersion
which induces an isomorphism of $Y$ with the inverse image
of $D(g) \subset \Spec(A_0)$. Moreover the map $\theta$
is an isomorphism.
\end{lemma}

\begin{proof}
This is a special case of Lemma \ref{lemma-base-change-map-proj} above.
\end{proof}

\begin{lemma}
\label{lemma-d-uple}
Let $S$ be a graded ring. Let $d \geq 1$. Set $S' = S^{(d)}$ with notation
as in Algebra, Section \ref{algebra-section-graded}. Set
$X = \text{Proj}(S)$ and $X' = \text{Proj}(S')$. There is a canonical
isomorphism $i : X \to X'$ of schemes such that
\begin{enumerate}
\item for any graded $S$-module $M$ setting $M' = M^{(d)}$,
we have a canonical isomorphism $\widetilde{M} \to i^*\widetilde{M'}$,
\item we have canonical isomorphisms
$\mathcal{O}_{X}(nd) \to i^*\mathcal{O}_{X'}(n)$
\end{enumerate}
and these isomorphisms are compatible with the multiplication maps
of Lemma \ref{lemma-widetilde-tensor} and hence with the maps
(\ref{equation-multiply}),
(\ref{equation-multiply-on-sheaf}),
(\ref{equation-global-sections}),
(\ref{equation-global-sections-module}),
(\ref{equation-multiply-more-generally}), and
(\ref{equation-global-sections-more-generally}) (see proof for precise
statements.
\end{lemma}

\begin{proof}
The injective ring map $S' \to S$ (which is not a homomorphism of graded rings
due to our conventions), induces a map $j : \Spec(S) \to \Spec(S')$.
Given a graded prime ideal $\mathfrak p \subset S$ we see that
$\mathfrak p' = j(\mathfrak p) = S' \cap \mathfrak p$
is a graded prime ideal of $S'$.
Moreover, if $f \in S_+$ is homogeneous and $f \not \in \mathfrak p$, then
$f^d \in S'_+$ and $f^d \not \in \mathfrak p'$. Conversely, if
$\mathfrak p' \subset S'$ is a graded prime ideal not containing some
homogeneous element $f \in S'_+$, then
$\mathfrak p = \{g \in S \mid g^d \in \mathfrak p'\}$ is a
graded prime ideal of $S$ not containing $f$ whose image under $j$
is $\mathfrak p'$. To see that $\mathfrak p$ is an ideal, note
that if $g, h \in \mathfrak p$, then
$(g + h)^{2d} \in \mathfrak p'$ by the binomial formula
and hence $g + h \in \mathfrak p'$ as $\mathfrak p'$ is a prime.
In this way we see that $j$ induces a homeomorphism $i : X \to X'$.
Moreover, given $f \in S_+$ homogeneous, then we have
$S_{(f)} \cong S'_{(f^d)}$. Since these isomorphisms are compatible
with the restrictions mappings of
Lemma \ref{lemma-standard-open}, we see that there exists an
isomorphism $i^\sharp : i^{-1}\mathcal{O}_{X'} \to \mathcal{O}_X$ of
structure sheaves on $X$ and $X'$, hence $i$ is an isomorphism
of schemes.

\medskip\noindent
Let $M$ be a graded $S$-module. Given $f \in S_+$ homogeneous, we have
$M_{(f)} \cong M'_{(f^d)}$, hence in exactly the same manner as above
we obtain the isomorphism in (1). The isomorphisms in (2) are a special
case of (1) for $M = S(nd)$ which gives $M' = S'(n)$. Let $M$ and $N$
be graded $S$-modules. Then we have
$$
M' \otimes_{S'} N' =
(M \otimes_S N)^{(d)} =
(M \otimes_S N)'
$$
as can be verified directly from the definitions. Having said this
the compatibility with the multiplication maps of
Lemma \ref{lemma-widetilde-tensor} is the commutativity of the diagram
$$
\xymatrix{
\widetilde M \otimes_{\mathcal{O}_X} \widetilde N
\ar[d]_{(1) \otimes (1)} \ar[r] &
\widetilde{M \otimes_S N} \ar[d]^{(1)} \\
i^*(\widetilde{M'} \otimes_{\mathcal{O}_{X'}} \widetilde{N'}) \ar[r] &
i^*(\widetilde{M' \otimes_{S'} N'})
}
$$
This can be seen by looking at the construction of the maps
over the open $D_+(f) = D_+(f^d)$ where the top horizontal
arrow is given by the map
$M_{(f)} \times N_{(f)} \to (M \otimes_S N)_{(f)}$
and the lower horizontal arrow by the map
$M'_{(f^d)} \times N'_{(f^d)} \to (M' \otimes_{S'} N')_{(f^d)}$.
Since these maps agree via the identifications
$M_{(f)} = M'_{(f^d)}$, etc, we get the desired compatibility.
We omit the proof of the other compatibilities.
\end{proof}











